\section{附注}\label{sec:notes-on-riemannian-manifolds}

\subsection{第 \ref{sec:hopf-rinow-theorem} 节}\label{subsec:notes-hopf-rinow}

本节用到几个关于度量空间的定义和结果, 现在来讨论它们. \textbf{度量空间}是一个二元组 $(M,d)$, 其中 $M$ 是一个集合, $d:M\times M\to[0,+\infty)$ 是满足 Proposition~\ref{prop:riemannian-distance-defines-metric-topology} 所列性质的映射. 集合
\[
  B_\varepsilon(p)=\{q\in M\mid d(p,q)<\varepsilon\}
\]
称为以 $p$ 为中心, 半径为 $\varepsilon$ 的\textbf{开球}. 所有这类球的族构成 $M$ 上一个 Hausdorff 拓扑的基, 此拓扑称为\textbf{度量拓扑}. 注意, 在此拓扑中, $p_n\to p$ 当且仅当 $d(p_n,p)\to0$. 虽然度量空间 $(M,d)$ 不一定满足第二可数性, 下述结论仍然成立: $F\subset M$ 为闭集当且仅当 $F$ 中的每个收敛序列的极限都在 $F$ 中, 而 $K\subset M$ 紧当且仅当 $K$ 中的每个序列在 $K$ 中都有子列极限.

若对所有 $\varepsilon>0$, 存在 $N\in\mathbb{N}$, 使得对所有 $m,n>N$ 都有 $d(p_n,p_m)<\varepsilon$, 则称 $M$ 中的序列 $\{p_n\}$ 为 \textbf{Cauchy 序列}. 容易看出所有收敛序列都是 Cauchy 序列. 然而反过来不一定成立 (但若一个 Cauchy 序列有收敛子序列, 则它必收敛). 若度量空间的所有 Cauchy 序列均收敛, 就称它\textbf{完备}. 完备度量空间的闭子集本身也是完备的.

若集合是某个球的子集, 就称它\textbf{有界}. 例如, Cauchy 序列所有项所成的集合有界. 容易证明, 若 $K\subset M$ 紧, 则 $K$ 必有界且闭 (但反过来不一定成立). 紧度量空间必完备.

\subsection{文献附注}\label{subsec:riemannian-manifolds-literature}

本章的材料在大多数 Riemann 几何著作中都可以找到 (例如 [\RGref{bib:ch03-Boo03}, \RGref{bib:ch03-dC93}, \RGref{bib:ch03-GHL04}]). 一般仿射联络的更多细节参见 \RGcite{KN96}. Lie 群上的双不变度量是对称空间的例子, 其优美的理论例如可见 \RGcite{Hel01}.
